
\documentclass[12pt]{modelo/prova2}

\orientacoes{
1. Leia cada questão com atenção. \\
2. Marque apenas uma resposta correta com um X. \\
3. Se precisar de ajuda para entender a pergunta, sinalize para a professora. \\[.5em] 
\textbf{Boa prova!}
}

\begin{document}
\maketitle

\question[1.0] O endereço \textbf{IPv4} é a identificação de um dispositivo na rede. Quantos \textbf{bits} existem em um endereço IPv4?
\begin{choices}
    \item 16 bits.
    \item 32 bits.
    \item 64 bits.
    \item 128 bits.
\end{choices}
\answer{B}

\question[1.0] Os protocolos de roteamento usam \textbf{métricas} para escolher o melhor caminho. Relacione o protocolo à sua métrica correta:

\begin{center}
\begin{tabular}{ |p{2cm}|p{6cm}| } 
 \hline
 \textbf{Protocolo}& \textbf{Métrica} \\ 
 \hline
 RIP  & Contagem de Saltos (Hops) \\ 
 OSPF & Custo (Baseado em Banda) \\ 
 \hline
\end{tabular}
\end{center}
Com base na tabela acima, qual protocolo escolhe o caminho com o \textbf{menor número de roteadores no meio}?

\begin{choices}
\item Protocolo IP.
\item Protocolo OSPF.
\item Protocolo RIP.
\item Protocolo DHCP.
\end{choices}
\answer{C}

\question[1.0] O roteador usa o \textbf{Prefix Matching} (Casamento de Prefixo) para enviar o pacote. Observe a tabela simplificada abaixo:
\begin{center}
\begin{tabular}{ |p{4cm}|p{4cm}| } 
 \hline
 \textbf{Prefixo de Destino}& \textbf{Interface de Saída}\\ 
 \hline
 11001000 & Interface 0 \\
 11001001 & Interface 1 \\
 \hline
\end{tabular}
\end{center}
Se um pacote chegar com o destino \verb|11001001|, para qual interface ele será enviado?
\begin{choices}
\item Interface 0
\item Interface 1
\item Nenhuma interface
\end{choices}
\answer{B}
\newpage
\question[1.0]  O \textbf{IPv6} foi criado porque os endereços IPv4 acabaram. Qual é a principal diferença do IPv6 em relação ao IPv4?
\begin{choices}
\item O IPv6 possui muito mais endereços (128 bits) e não usa mais o Broadcast.
\item O IPv6 é menor que o IPv4.
\item O IPv6 só funciona em redes sem fio.
\item O IPv6 é mais lento que o IPv4.
\end{choices}
\answer{B}

\question[1.0] O protocolo \textbf{DHCP} serve para entregar endereços IP automaticamente. Qual é a \textbf{ordem correta} das mensagens enviadas?
\begin{choices}
    \item Discover $\to$ Offer $\to$ Request $\to$ Ack.
    \item Request $\to$ Offer $\to$ Discover $\to$ Ack.
    \item Ack $\to$ Request $\to$ Offer $\to$ Discover.
    \item Offer $\to$ Discover $\to$ Request $\to$ Ack.
\end{choices}
\answer{A}

\question[1.0] Nas redes \textbf{SDN} (Redes Definidas por Software), dividimos as funções do roteador. Qual plano é responsável por \textbf{decidir para onde o pacote deve ir} (o "cérebro" da rede)?
\begin{choices}
\item Plano de Dados (Data Plane).
\item Plano de Controle (Control Plane).
\item Plano de Energia.
\item Plano de Aplicação.
\end{choices}
\answer{B}

\question[1.0] O protocolo \textbf{BGP} é usado para conectar grandes redes na Internet. Como o BGP escolhe o caminho?
\begin{choices}
    \item Ele escolhe sempre o caminho mais curto em termos de saltos (hops).
    \item Ele escolhe sempre o caminho mais rápido, como um GPS.
    \item Ele escolhe o caminho de forma aleatória.
    \item Ele escolhe baseado em políticas comerciais e acordos entre empresas. 
\end{choices}
\answer{D}

\question[1.0] Você precisa configurar uma rede para um departamento com 60 computadores (hosts). Qual máscara de sub-rede (VLSM) é a mais adequada para suportar exatamente essa quantidade?
\begin{choices}
    \item /25 (Suporta até 126 hosts).
    \item /26 (Suporta até 62 hosts).
    \item /27 (Suporta até 30 hosts).
    \item /28 (Suporta até 14 hosts).
\end{choices}
\answer{B}

\question[1.0] O protocolo \textbf{ICMP} é usado para enviar mensagens de erro e teste. Qual comando usamos para \textbf{testar se um computador está conectado} na rede?
\begin{choices}
\item Comando DHCP.
\item Comando PING.
\item Comando Traceroute.
\item Comando FTP.
\end{choices}
\answer{B}

\question[1.0] No Modelo OSI e no Modelo TCP/IP, o \textbf{Roteador} é um equipamento principal. Em qual camada o roteador trabalha?
\begin{choices}
\item Camada de Aplicação.
\item Camada de Transporte.
\item Camada de Rede (ou Camada Internet).
\item Camada Física.
\end{choices}
\answer{C}


\end{document}
